\ResearchStyle
\PartDivider{PART I}{Research Background}{Build the networking foundation needed to understand CATRA.}

\begin{frame}[shrink=4]{Multi-Hop 802.11 Ad Hoc Networks}

\KeyPoint{In ad hoc mode, stations communicate without an access point and relay packets over multiple wireless hops.}

\vspace{0.18cm}

\begin{columns}[T,onlytextwidth]
\column{0.55\textwidth}
\begin{center}
\begin{tikzpicture}[
    station/.style={circle,draw=ResearchColor,fill=CardBg,minimum size=0.82cm,font=\bfseries},
    link/.style={draw=ResearchColor,line width=1pt},
    route/.style={-{Latex[length=2.2mm]},draw=WarningColor,line width=1.5pt}
]
    \node[station] (a) at (0,0) {A};
    \node[station] (b) at (1.65,0.65) {B};
    \node[station] (c) at (3.30,0) {C};
    \node[station] (d) at (4.95,0.65) {D};

    \draw[link] (a) -- (b);
    \draw[link] (b) -- (c);
    \draw[link] (c) -- (d);
    \draw[route] ([yshift=-0.18cm]a.east) -- ([yshift=-0.18cm]b.west);
    \draw[route] ([yshift=-0.18cm]b.east) -- ([yshift=-0.18cm]c.west);
    \draw[route] ([yshift=-0.18cm]c.east) -- ([yshift=-0.18cm]d.west);

    \node[font=\fontsize{8}{9}\selectfont,text=SecondaryColor,align=center] at (2.48,-0.95)
    {One end-to-end packet requires\\three channel-access events};
\end{tikzpicture}
\end{center}

\column{0.41\textwidth}
\fontsize{9}{10.8}\selectfont
\textbf{Key characteristics}
\begin{itemize}
    \item No centralized scheduler or fixed infrastructure.
    \item Nodes act as both end hosts and packet forwarders.
    \item Every hop consumes airtime in a shared broadcast medium.
    \item Neighboring hops cannot always transmit simultaneously.
\end{itemize}
\end{columns}

\vspace{0.12cm}
\SmallNote{\textbf{Implication:} End-to-end throughput depends on route length and local contention around every forwarding node.}

\end{frame}


\begin{frame}[shrink=3]{Single-Hop vs. Multi-Hop Communication}
\begin{columns}[T,onlytextwidth]
\column{0.45\textwidth}
\textbf{Single-hop}
\begin{center}
\begin{tikzpicture}
\node[draw=ResearchColor,circle,minimum size=0.8cm] (s) {S};
\node[draw=ResearchColor,circle,minimum size=0.8cm,right=3.0cm of s] (r) {R};
\draw[-{Latex[length=2.5mm]},draw=ResearchColor,line width=1.2pt] (s)--(r);
\end{tikzpicture}
\end{center}
\begin{itemize}
\item One packet needs one channel-access event.
\item Sender and receiver share one local contention domain.
\end{itemize}

\column{0.51\textwidth}
\textbf{Multi-hop}
\begin{center}
\begin{tikzpicture}[node distance=0.72cm]
\node[draw=ResearchColor,circle,minimum size=0.72cm] (s) {S};
\node[draw=ResearchColor,circle,minimum size=0.72cm,right=of s] (n1) {N1};
\node[draw=ResearchColor,circle,minimum size=0.72cm,right=of n1] (n2) {N2};
\node[draw=ResearchColor,circle,minimum size=0.72cm,right=of n2] (r) {R};
\draw[-{Latex[length=2mm]},draw=ResearchColor] (s)--(n1);
\draw[-{Latex[length=2mm]},draw=ResearchColor] (n1)--(n2);
\draw[-{Latex[length=2mm]},draw=ResearchColor] (n2)--(r);
\end{tikzpicture}
\end{center}
\begin{itemize}
\item The same packet consumes airtime at every forwarding hop.
\item Neighboring hops contend and often cannot transmit together.
\end{itemize}
\end{columns}
\vspace{0.35cm}
\KeyPoint{Route length affects airtime consumption, contention, and end-to-end throughput.}
\end{frame}

\begin{frame}{Why a Shared Wireless Medium Is Different}
\begin{columns}[T,onlytextwidth]
\column{0.48\textwidth}
\textbf{Broadcast and half-duplex}
\begin{itemize}
\item A transmission can be sensed by multiple neighboring stations.
\item A radio normally cannot detect a collision while transmitting.
\item Wi-Fi therefore avoids collisions instead of detecting them.
\end{itemize}
\column{0.48\textwidth}
\textbf{Multi-hop consequences}
\begin{itemize}
\item Forwarding traffic competes with locally generated traffic.
\item Spatial reuse depends on physical and carrier-sensing ranges.
\item PHY rate is not equal to end-to-end goodput.
\end{itemize}
\end{columns}
\vspace{0.35cm}
\begin{center}
\begin{tikzpicture}[node distance=0.65cm]
\node[draw=ResearchColor,rounded corners,fill=CardBg] (a) {More hops};
\node[draw=ResearchColor,rounded corners,fill=CardBg,right=of a] (b) {More channel uses};
\node[draw=WarningColor,rounded corners,fill=CardBg,right=of b] (c) {More contention};
\node[draw=WarningColor,rounded corners,fill=CardBg,right=of c] (d) {Lower E2E rate};
\draw[-{Latex[length=2mm]},draw=ResearchColor] (a)--(b);
\draw[-{Latex[length=2mm]},draw=WarningColor] (b)--(c);
\draw[-{Latex[length=2mm]},draw=WarningColor] (c)--(d);
\end{tikzpicture}
\end{center}
\end{frame}

\begin{frame}[shrink=3]{IEEE 802.11 MAC: Distributed Channel Access}

\KeyPoint{DCF uses carrier sensing, random backoff, and acknowledgements to coordinate access without centralized scheduling.}

\vspace{0.05cm}

\begin{center}
\begin{tikzpicture}[
    node distance=0.52cm,
    macstep/.style={
        draw=ResearchColor,
        rounded corners=3pt,
        fill=CardBg,
        minimum width=2.25cm,
        minimum height=0.78cm,
        align=center,
        font=\fontsize{8}{9.5}\selectfont\bfseries
    },
    flow/.style={-{Latex[length=2.2mm]},draw=ResearchColor,line width=0.75pt}
]
    \node[macstep] (sense) {Sense the\\channel};
    \node[macstep,right=of sense] (wait) {Wait for\\DIFS};
    \node[macstep,right=of wait] (backoff) {Random\\backoff};
    \node[macstep,right=of backoff] (send) {Transmit\\frame};
    \node[macstep,right=of send] (ack) {Receive\\ACK};

    \draw[flow] (sense) -- (wait);
    \draw[flow] (wait) -- (backoff);
    \draw[flow] (backoff) -- (send);
    \draw[flow] (send) -- (ack);
\end{tikzpicture}
\end{center}

\vspace{0.12cm}

\begin{columns}[T,onlytextwidth]
\column{0.48\textwidth}
\fontsize{9}{10.8}\selectfont
\textbf{CSMA/CA and backoff}
\begin{itemize}
    \item A station selects a random counter from its contention window (CW).
    \item The counter decreases in idle slots and freezes whenever the channel becomes busy.
\end{itemize}

\column{0.48\textwidth}
\fontsize{9}{10.8}\selectfont
\textbf{When contention increases}
\begin{itemize}
    \item The receiver returns an ACK after SIFS; a missing ACK indicates loss or collision.
    \item Binary Exponential Backoff enlarges CW before retransmission.
\end{itemize}
\end{columns}

\end{frame}


\begin{frame}[shrink=5]{Contention Window and Expected Backoff}
\KeyPoint{A station draws an integer backoff counter from its current contention window.}
\[
N\sim U\{0,\ldots,CW\},\qquad
E[T_{backoff}]\approx \frac{CW}{2}\,SlotTime
\]
\begin{columns}[T,onlytextwidth]
\column{0.48\textwidth}
\textbf{Smaller CW}
\begin{itemize}
\item Shorter expected wait.
\item More aggressive channel access.
\item Higher collision risk under contention.
\end{itemize}
\column{0.48\textwidth}
\textbf{Larger CW}
\begin{itemize}
\item Longer expected wait.
\item Lower access probability.
\item More room for competing stations.
\end{itemize}
\end{columns}
\vspace{0.25cm}
\KeyPoint{CATRA changes CW deliberately to control each station's channel-access opportunity.}
\end{frame}

\begin{frame}{Binary Exponential Backoff}
\begin{center}
\begin{tikzpicture}[node distance=0.55cm]
\node[draw=ResearchColor,rounded corners,fill=CardBg] (tx) {Transmit};
\node[draw=WarningColor,rounded corners,fill=CardBg,right=of tx] (fail) {No ACK};
\node[draw=WarningColor,rounded corners,fill=CardBg,right=of fail] (cw) {$CW\leftarrow\min(2CW,CW_{max})$};
\node[draw=ResearchColor,rounded corners,fill=CardBg,right=of cw] (retry) {New backoff};
\draw[-{Latex[length=2mm]}] (tx)--(fail);
\draw[-{Latex[length=2mm]}] (fail)--(cw);
\draw[-{Latex[length=2mm]}] (cw)--(retry);
\end{tikzpicture}
\end{center}
\vspace{0.35cm}
\begin{itemize}
\item BEB reacts to success/failure and contention conditions.
\item It does not explicitly know the number of neighboring flows.
\item It does not know whether a forwarding station receives its fair share.
\item Consequently, the resulting CW may be inappropriate for fairness in asymmetric multi-hop topologies.
\end{itemize}
\end{frame}

\begin{frame}{802.11 Timing Used by CATRA}
\begin{center}
\begin{tikzpicture}[x=1.0cm,y=1cm]
\draw[-{Latex[length=2mm]},draw=SecondaryColor] (0,0)--(12.3,0);
\foreach \x/\w/\label/\c in {0.2/1.2/DIFS/ResearchColor,1.4/2.0/Backoff/ResearchColor,3.4/1.1/RTS/WarningColor,4.5/0.7/SIFS/ResearchColor,5.2/1.1/CTS/WarningColor,6.3/0.7/SIFS/ResearchColor,7.0/2.4/TCP DATA/TitleColor,9.4/0.7/SIFS/ResearchColor,10.1/1.2/MAC ACK/WarningColor}
{\fill[\c!18] (\x,0.12) rectangle +(\w,0.62);\draw[\c] (\x,0.12) rectangle +(\w,0.62);\node[font=\fontsize{7}{8}\selectfont] at (\x+\w/2,0.43) {\label};}
\end{tikzpicture}
\end{center}
\vspace{0.45cm}
\begin{itemize}
\item Algorithm 1 approximates airtime from average backoff, control frames, inter-frame spaces, TCP DATA/ACK, and MAC ACK.
\item The receiver's TCP ACK travels in the reverse direction and also consumes wireless airtime.
\item Active time is therefore a resource-usage estimate, not only a packet counter.
\end{itemize}
\end{frame}

\input{slides/10-wireless-contention}
\input{slides/10-transport-layer}
\begin{frame}[shrink=9]{TCP Tahoe Congestion Control}

\KeyPoint{TCP Tahoe uses packet loss as a congestion signal and controls the sender through $cwnd$ and $ssthresh$.}

\vspace{0.12cm}

\begin{columns}[T,onlytextwidth]
\column{0.56\textwidth}
\begin{center}
\begin{tikzpicture}[x=0.52cm,y=0.115cm]
    \draw[-{Latex[length=2mm]},draw=SecondaryColor] (0,0) -- (10.5,0)
        node[right,font=\fontsize{7}{8}\selectfont] {RTT};
    \draw[-{Latex[length=2mm]},draw=SecondaryColor] (0,0) -- (0,31)
        node[above,font=\fontsize{7}{8}\selectfont] {$cwnd$};
    \draw[draw=ResearchColor,line width=1.5pt]
        (0.4,2) -- (1.4,4) -- (2.4,8) -- (3.4,16) -- (4.4,24)
        -- (5.4,27) -- (6.4,30);
    \draw[draw=WarningColor,line width=1.5pt,dashed] (6.4,30) -- (6.7,2);
    \draw[draw=ResearchColor,line width=1.5pt]
        (6.7,2) -- (7.7,4) -- (8.7,8) -- (9.7,15);
    \draw[draw=SecondaryColor,dashed] (0,16) -- (10,16);
    \node[font=\fontsize{7.5}{8.5}\selectfont,text=ResearchColor] at (2.4,21) {Slow start};
    \node[font=\fontsize{7.5}{8.5}\selectfont,text=ResearchColor] at (5.4,25) {Additive increase};
    \node[font=\fontsize{7.5}{8.5}\selectfont,text=WarningColor,align=center] at (7.7,26) {Loss: halve $ssthresh$,\\reset $cwnd$};
    \node[font=\fontsize{7}{8}\selectfont,text=SecondaryColor] at (1.4,17.5) {$ssthresh$};
\end{tikzpicture}
\end{center}

\column{0.40\textwidth}
\fontsize{9}{10.8}\selectfont
\textbf{Control phases}
\begin{itemize}
    \item \textbf{Slow start:} $cwnd$ grows exponentially with ACKs.
    \item \textbf{Congestion avoidance:} growth becomes approximately one segment per RTT.
    \item \textbf{Loss response:} set $ssthresh \leftarrow cwnd/2$ and restart from a small window.
    \item Three duplicate ACKs trigger fast retransmission.
\end{itemize}
\end{columns}

\end{frame}


\begin{frame}{Station Fairness vs. Flow Fairness}
\begin{columns}[T,onlytextwidth]
\column{0.48\textwidth}
\textbf{Per-station fairness}
\begin{center}
\begin{tikzpicture}
\node[draw=ResearchColor,rounded corners,fill=CardBg,text width=4.8cm,align=center,minimum height=1.2cm]
{Does each contending station receive a fair channel-access opportunity?};
\end{tikzpicture}
\end{center}
\column{0.48\textwidth}
\textbf{Per-flow fairness}
\begin{center}
\begin{tikzpicture}
\node[draw=ResearchColor,rounded corners,fill=CardBg,text width=4.8cm,align=center,minimum height=1.2cm]
{Does each end-to-end flow receive a fair throughput share?};
\end{tikzpicture}
\end{center}
\end{columns}
\vspace{0.45cm}
\KeyPoint{One station may forward several flows; therefore equal station airtime does not imply equal flow throughput.}
\end{frame}

\begin{frame}[shrink=5]{Why MAC and TCP Must Be Considered Together}

\KeyPoint{In a multi-hop wireless path, transport-layer load changes MAC contention, while MAC behavior changes TCP feedback.}

\vspace{0.25cm}

\begin{center}
\begin{tikzpicture}[
    node distance=0.62cm,
    stage/.style={draw=ResearchColor,rounded corners=4pt,fill=CardBg,text width=2.55cm,
        minimum height=0.92cm,align=center,font=\fontsize{8.5}{10}\selectfont\bfseries},
    effect/.style={draw=WarningColor,rounded corners=4pt,fill=CardBg,text width=2.55cm,
        minimum height=0.92cm,align=center,font=\fontsize{8.5}{10}\selectfont\bfseries},
    flow/.style={-{Latex[length=2.2mm]},line width=0.9pt}
]
    \node[stage] (cwnd) {Larger TCP\\congestion window};
    \node[stage,right=of cwnd] (queue) {More packets at\\each wireless hop};
    \node[effect,right=of queue] (contention) {More contention,\\backoff, and retries};
    \node[effect,below=0.85cm of contention] (feedback) {Higher RTT and\\packet-loss signals};
    \node[stage,left=of feedback] (tcp) {TCP reduces or\\regrows $cwnd$};

    \draw[flow,draw=ResearchColor] (cwnd) -- (queue);
    \draw[flow,draw=WarningColor] (queue) -- (contention);
    \draw[flow,draw=WarningColor] (contention) -- (feedback);
    \draw[flow,draw=ResearchColor] (feedback) -- (tcp);
    \draw[flow,draw=ResearchColor] (tcp.west) -| ([xshift=-0.35cm]cwnd.south) -- (cwnd.south);
\end{tikzpicture}
\end{center}

\vspace{0.25cm}
\begin{center}
\fontsize{9}{10.8}\selectfont
\textbf{Required coordination:} fair station-level channel access at MAC + fair flow-level demand at TCP.
\end{center}

\end{frame}

